\subsection{Structure Embedding}\label{methods/feature_engineering/structure_embedding}

To represent the structure of a given protein the \gls{prostT5} presented by \citet{Heinzinger2023}, model was used. \gls{prostT5} is a version of the ProtT5 model finetuned to adapt to the task of bidirectional translation of sequence and structure. \autoref{fig:prostT5Architecture} illustrates the overall pipeline of the protein language model. \gls{prostT5} is meant to enable direct translation of the amino acid sequence of a protein and it's structure. The structure here is represented using the 3Di a structure alphabet developed by Foldseek to enable efficient structural humolgy search \citep{Kempen2024}. 

A target protein structure is discretized into a sequence over the 3Di structural alphabet. Each token in the 3Di sequence represents the local tertiary interaction geometry of a residue with its nearest neighbor in three-dimensional space. The alphabet consists of 20 discrete states that were optimized to maximize structural information content and search sensitivity.

A target protein structure is discretized into a sequence over the 3Di structural alphabet. Each token in the 3Di sequence represents the local \emph{tertiary interaction geometry} of a residue with its nearest neighbor in three-dimensional space. The alphabet consists of 20 discrete states that were optimized to maximize structural information content and search sensitivity.

For every residue $i$ in a structure, the nearest spatial neighbor $j$ is identified based on a residue-specific virtual center. The local interaction between residues $i$ and $j$ is then described using a 10-dimensional feature vector consisting of angular features, the Euclidean distance between residues, and sequence-distance features. This continuous descriptor vector is subsequently discretized into one of the 20 3Di states using a trained encoder.

In particular \citet{Kempen2024}, the 3Di alphabet was not manually designed but learned using a vector-quantized variational autoencoder (VQ-VAE). The network consists of an encoder and a decoder, with a discrete bottleneck in between representing the latent 3Di space. Each of the 20 characters in the 3Di alphabet corresponds to a learned centroid in this latent space. During discretization, a new descriptor vector is mapped to the nearest centroid, yielding its corresponding 3Di symbol.

Because continuous geometric information is quantized into a finite set of discrete states, the representation is inherently lossy. Nevertheless, the authors demonstrate that structure homology search based on 3Di sequences achieves sensitivity comparable to state-of-the-art structural alignment methods, while being orders of magnitude $faster$.

\begin{figure}[htbp]
    \centering
    \includegraphics[width=0.8\linewidth]{media/ProstT5.png}
    \caption{Description of ProstT5 pipeline}{The figure visualizes the training, architectre and inference pipeline of \gls{prostT5}.}
    \label{fig:prostT5Architecture}
\end{figure}

Fig. 1: Sketch of ProstT5. Model architecture: ProstT5 is a T5-based encoder-decoder model initialized with weights of the
ProtT5 model6. Pre-training: Foldseek1 transferred protein 3D coordinates into 3Di tokens, i.e., 1D descriptions of 3D structure
that assign each residue in a protein into one of twenty states described by a 1D string of letters. We used 17 million (17M)
high-quality, non-redundant and diverse 3D predictions from AFDB33. ProtT5 was leveraged as an already pre-trained starting
point for translating between 1D sequence (amino acids, AA) and 3D structure (3Di). Firstly, we applied the original pre-training
objective of ProtT5 (span-based denoising) to both, AAs and 3Di, to teach the model the new 3Di tokens while avoiding
catastrophic forgetting of AAs. Secondly, we continued to train the resulting model to translate between AAs and 3Di and vice
versa. The final model, ProstT5 (Protein structure-sequence T5) extracts the information in its internal embeddings that can be
input into downstream applications. This includes established feature extraction using only the encoder6, or bi-directional
translation, either from AAs to 3Di (“folding”) or from 3Di to AAs (“inverse folding”). Inference: bi-directional translation from AA
to 3Di (AA→3Di) or 3Di→AA can be conducted using either the encoder-decoder mode, necessitating token-wise
decoder-inference or through an optimized inference mode, where 3Di tokens are directly predicted through a convolutional
neural network from the encoder-embedding. The optimized 3Di inference mode results in a three orders of magnitude
speedup over 3Di extraction from predicted protein structures (Fig. 2).